\section{引言}
\label{sec:intro}

再生数是判断传染病传播态势的基本指标。它大于 1 意味着疫情扩张，小于 1 意味着收缩，其变化常被用来评估干预措施的效果和安排医疗资源。基于世代间隔和更新方程的估计方法已成为实时追踪的主要工具\citep{wallinga2004different, fraser2007estimating, cori2013new}，窗口选择、报告延迟、输入病例和偏差修正等实际问题也有系统讨论\citep{thompson2019improved, gostic2020practical, nouvellet2018simple, steyn2025primer}。这些工作大多回答给定一条病例序列后怎样估计得更好，监测系统设计者面对的是前一步的问题，即应当收集多少病例、覆盖多少机构，才能让估计达到所需的精度。

规模设计在住院监测中尤其具体。流感住院只占感染的很小一部分，按美国的估计，住院比例约为 0.6\%--2\%\citep{zhou2012hospitalizations, tokars2018seasonal}，每个辖区每周的住院数往往只有几十例。美国国家医疗安全网络（National Healthcare Safety Network，NHSN）的流感住院周报在 2024 年经历了从强制填报到自愿填报、再恢复强制填报的切换\citep{cdc2026nhsn}，填报医院比例随之大幅波动。以这些数据为输入的集合预测项目已形成系统的评价框架\citep{reich2019accuracy, cramer2022evaluation, mathis2024evaluation}，相对误差类评分的性质也有专门研究\citep{bosse2023scoring}。监测部门因此需要知道，误差中有多大部分能靠扩大规模消除，提高报告率与扩大覆盖面哪一条途径更划算，规模扩大到何处以后收益明显下降。

经典的统计直觉给出了一个简单答案。在泊松更新模型下，再生数的 Fisher 信息随病例数线性增长，估计方差与病例数成反比\citep{parag2022quantifying}。超传播使继发感染数严重过离散，通常用负二项分布刻画\citep{lloydsmith2005superspreading, endo2020estimating, kucharski2020early}；它改变方差的比例常数，方差随病例数下降的速度保持不变。按照这一直觉，扩大监测规模总能换来同比例的精度提升，规模的边际收益只随规模平滑递减，不存在由传播过程本身决定的转折点。

真实监测数据中还有两类噪声，它们使上述直觉需要修正。第一类是环境扰动。同一辖区、同一时段的感染者共享天气、接触模式和人群流动。学校开学、假期以及湿度与气温的变化使流感传播呈现明显的季节性，各期再生数因而围绕均值波动\citep{dalziel2018urbanization, rasmussen2011inference}。生态学把这类波动称为环境随机性，并与个体层面的人口随机性相区分。人口随机性的相对影响随种群规模增大而消失，环境随机性则不随规模减弱\citep{engen1998demographic, rohani2002determinism}。病例数增加能把瞬时再生数测得更准，却无法消除再生数本身的跨期波动。若监测的目的是把握一段时期的典型传播水平，或预测下一期的传播，这部分波动就成为扩大规模无法消除的误差。

第二类噪声来自数据结构。下一周的感染压力取决于此前各周的感染人数，漏报使这些人数无法完全观测，估计所用的分母因而也是随机的。针对报告噪声，已有工作发展了漏报和报告偏差的修正方法\citep{azmon2014estimating, stoner2019hierarchical, bracher2021marginal}，以及状态空间模型与半机制模型\citep{breto2009time, bhatt2023semi, flaxman2020estimating}；Parag 等\citep{parag2021improved}的 EpiFilter 用相邻时期的信息平滑 $R_t$ 估计。Parag 等\citep{parag2022quantifying, parag2022fundamental, parag2025asymmetric}量化了带噪发病曲线所含的信息，并给出识别疫情反弹所需时间的下限。上述工作关心给定数据能提供多少关于当期传播的信息，环境波动没有作为单独的方差成分出现，规模扩大到何种程度后边际收益明显下降，也就没有定量回答。分母噪声如何影响以病例数为横轴的误差曲线，同样缺少可以用数据检验的预测。报告比例本身也随就诊行为、检测能力和填报制度逐期变化，这部分波动作用在整期报告数上，扩大覆盖人数同样无法消除，它与环境方差在计数中如何区分，也没有现成的答案。

针对上述缺口，本文以负二项传播和二项报告为基本模型，分别处理两个估计目标。估计目标1是给定当期环境时实际发生的瞬时再生数 $R_t$；估计目标2是各期再生数所围绕的跨期平均再生数 $\Rbar$。估计目标2与一步前瞻预测等价有两个条件：各期环境独立（$\phi=0$），并且 $\Rbar$ 已知；以此前 $T_0$ 期的估计代替时误差乘以 $1+1/T_0$，环境存在时间相关时一步预测的误差下界降为 $(1-\phi^2)v_R$（第~\ref{sec:meth-targets}~节）。对估计目标1，$\widehat R_t=C_t/m$ 达到 Cramér--Rao 下界，方差按 $1/m$ 趋于零，离散参数趋于无穷时退化为 Parag 等的泊松结果。对估计目标2，报告率固定时条件无偏估计的方差等于环境方差 $v_R$ 与规模项 $A/m$ 之和，报告率随期波动时多出同样不随规模下降的报告异质性方差 $v_H$，规模收益开始明显减小的位置由方差平衡规模 $\mx=A/v_R$ 给出。这些结论来自 Cramér--Rao 不等式、全方差公式、Fisher 信息的数据处理不等式与单调似然比检验\citep{casella2002statistical, lehmann2005testing, zamir1998proof}，预算分配与两阶段抽样的费用分配同构\citep{cochran1977sampling}。

报告率逐期波动时，提高报告率与扩大覆盖的优劣存在反转阈值：二项型波动下单期规模超过 $\rho^2\nu/k$ 后提高报告率占优，制度型波动下扩大覆盖始终占优，而两条途径都不改变 $v_H$（推论~\ref{cor:switch}）。仅凭计数的前两阶矩无法区分 $v_R$ 与 $v_H$，但报告噪声使预测误差的截距按函数 $F(w,\psi)$ 随滞后权重变化，这一签名可以用滞后权重不同的分母检验（命题~\ref{prop:report-noise}）。规模足够大以后，$v_R$ 本身的估计精度只由有效期数 $T_{\mathrm{eff}}\approx T(1-\phi^2)/(1+\phi^2)$ 决定，季层面方差的相对标准误不低于 $\sqrt{2/S}$（命题~\ref{prop:vR-info} 与推论~\ref{cor:season}）。环境正相关时，预算低于阈值 $\mathcal B^\dagger$ 应设更少、更大的观测期，高于阈值才应设更多、更小的观测期（命题~\ref{prop:ar-design}）。

应用部分使用美国国家医疗安全网络 16,218 个州周的流感住院周报。流行初期的单步预测误差随规模按 $a+b/D_t$ 的形式下降，经验比值 $b/a$ 为 22.6 例/周（95\% 百分位区间 13.5--34.9），与情景参数给出的方差平衡规模处于同一量级，不含底板的均值形状则明显更差。截距高于环境方差情景，用各州每周填报医院比例调整计数后，可由填报比例观测到的报告成分不到截距超出量的 1\%，填报比例以外的报告离散不在这一上界之内。截距比随滞后权重变化的方向在流感、COVID-19 与 RSV 的 7 个病原与窗口组合中一致，但入院延迟在环境各期独立时也会产生这一方向。按截距比与截距超出量的联合判据，4 个组合部分复现，3 个组合只呈方向一致，没有组合复现。季间截距差异显著，季层面方差的大小只能靠积累流行季确定。疫情预警与趋势估计对监测规模的需求不同，宜分开设计。
